\section{Conclusion}


In this paper, we present a unified dual-mask physical model to address the challenge of depth estimation in non-Lambertian light field scenes. Conventional epipolar plane image (EPI) (EPI) methods rely heavily on the Lambertian reflectance assumption, which inevitably fails when encountering specular, translucent, or scattering surfaces in real-world environments. To overcome this physical bottleneck, we introduce a comprehensive framework that explicitly decouples illumination, depth, and material properties, enabling robust depth inference across complex mixed scenes without relying on strict photo-consistency.

The specific contributions of this work are detailed as follows:
1) \textbf{Three-Layer Angular Signal Decomposition Physical Model}: We formulate a physics-driven decomposition model that decouples illumination, depth, and material information at both physical and geometric levels. By theoretically proving that geometric features possess higher discriminative power than spectral features under discrete angular resolutions, we establish the theoretical feasibility of extracting material characteristics from non-Lambertian surfaces.
2) \textbf{Adaptive Geometric Dual-Mask Routing Architecture}: We design a dual-mask routing mechanism that dynamically adapts to local surface reflectance properties. This architecture overcomes the degradation of traditional single-EPI models on non-Lambertian surfaces by processing Lambertian and non-Lambertian regions within a unified, adaptive pipeline, thereby enhancing estimation accuracy and robustness in complex mixed environments.
3) \textbf{Domain-Balanced Training and System-Level Validation}: We integrate a domain-balanced sampling strategy with the directional EPI backbone to mitigate data distribution disparities. Extensive evaluations demonstrate that the proposed unified framework achieves an overall mean absolute error (MAE) of 0.133 and an MAE of 0.081 in mixed urban scenarios, confirming its effectiveness in handling real-world scenes with diverse material compositions.

Future investigations will extend this physical decomposition paradigm to emerging 3D representation frameworks, such as neural radiance fields, to facilitate comprehensive scene understanding. Broadening the application of the adaptive dual-mask routing mechanism to other multi-view geometric tasks will also be explored to advance robust visual perception in unconstrained environments.


From a system-level perspective, the transition from continuous spectral analysis to discrete geometric decomposition addresses a fundamental mismatch between algorithmic assumptions and optical hardware constraints. Light field cameras, governed by the spatial-angular resolution trade-off, inherently capture sparse discrete angular grids (e.g., $9 \times 9$ or $5 \times 5$). Traditional frequency-domain approaches implicitly assume continuous angular manifolds, leading to severe spectral leakage and aliasing when applied to these constrained hardware setups. By bypassing the stringent Nyquist-Shannon requirements in the angular domain, our geometric formulation aligns the computational prior with the physical sampling limits of the optical system. Furthermore, replacing global Fourier transforms with localized geometric operations significantly reduces computational overhead, offering a highly efficient, hardware-friendly paradigm for real-time embedded light field processing.

Beyond architectural design, the adaptive dual-mask routing mechanism resolves a critical optimization bottleneck in mixed-scene depth estimation: gradient conflict. In unified networks, the loss gradients originating from Lambertian regions (driven by strict photo-consistency) and non-Lambertian regions (driven by complex scattering models) frequently point in opposing directions within the parameter space, resulting in sub-optimal convergence. By explicitly decoupling the feature manifolds via physically grounded masks, the routing strategy projects the optimization process into two distinct, non-interfering subspaces. This not only stabilizes training dynamics but also prevents the network from learning compromised, "averaged" representations that typically blur depth boundaries at material transitions. The mask generation acts as a dynamic capacity allocator, ensuring that specialized network branches are activated only when the local surface reflectance demands it, thereby maximizing parameter efficiency.

Despite these advancements, inherent physical ambiguities and sensor-level limitations remain. Defining a single "ground truth" depth for translucent and volumetric scattering materials is physically ill-posed, as photons interact with a continuous volume rather than a distinct geometric surface. While the residual layer successfully isolates the material signal, the depth assignment for such regions remains an approximation of the primary scattering interface. Additionally, the geometric decomposition relies on a minimum angular baseline to compute reliable linear and residual components; in extremely sparse configurations (e.g., $2 \times 2$ or $3 \times 3$ angular grids), the statistical robustness of angular correlation coherence inevitably degrades. Finally, saturated specular highlights often clip the sensor response, violating the linear physical relationship assumed in the signal decomposition. Future work will explore integrating temporal light field sequences to disambiguate volumetric depth and incorporating high dynamic range (HDR) priors to mitigate sensor saturation artifacts.

Despite the demonstrated efficacy of the unified dual-mask physical model in decoupling illumination, depth, and material properties, several limitations remain. First, the current three-layer angular decomposition assumes that the residual component $\epsilon(u,v)$ predominantly captures single-bounce non-Lambertian reflections. However, in scenes with complex inter-reflections or highly transparent media (e.g., thick glass or volumetric fog), light undergoes multiple scattering events. In such extreme cases, $\epsilon(u,v)$ becomes entangled with secondary geometric cues and depth-dependent caustics, potentially degrading the accuracy of the material classifier and the subsequent dual-mask routing. Second, while our geometric features outperform continuous spectral methods under standard $9 \times 9$ discrete sampling, their robustness diminishes under extremely sparse angular resolutions (e.g., $3 \times 3$ or $5 \times 5$). The discrete derivative operations become highly susceptible to sensor noise and quantization artifacts, leading to erroneous mask assignments. Additionally, the current material classifier relies heavily on the statistical distribution of synthetic datasets, which may exhibit a domain gap when applied to real-world light field captures with unknown sensor noise models, optical aberrations, and varying baseline configurations. Furthermore, from a systems perspective, performing pixel-wise physical decomposition and dual-branch inference introduces non-trivial computational overhead and memory bandwidth bottlenecks, restricting deployment in latency-critical applications such as real-time autonomous navigation or high-frame-rate robotic manipulation.

To address these challenges, future work will explore the integration of temporal dynamics and advanced 3D representations into the physical model. By extending the spatial-angular analysis to light field video sequences, temporal consistency can provide robust constraints to disentangle complex multi-bounce light paths and dynamic specularities from true surface geometry. We will also focus on unsupervised domain adaptation techniques to bridge the synthetic-to-real gap, ensuring the physical priors remain invariant across diverse light field camera architectures. Additionally, we plan to investigate the synergy between our discrete geometric decomposition and explicit 3D rendering frameworks, such as 3D Gaussian Splatting. Embedding the dual-mask routing mechanism directly into the differentiable rasterization pipeline could enable real-time, view-consistent depth estimation for dynamic non-Lambertian scenes while maintaining strict physical consistency. Finally, to alleviate the computational burden on edge devices, we aim to develop hardware-aware adaptive angular sampling strategies, where the network dynamically selects the minimal required angular views per spatial region based on local material complexity, significantly accelerating inference without compromising depth accuracy.